\documentclass[10pt,a4paper]{amsart}
\usepackage[margin=19mm]{geometry}
\usepackage{amsmath,amssymb,mathtools}
\usepackage[scr=rsfs,cal=euler]{mathalfa}
\usepackage{xfrac}
\usepackage[unicode,hidelinks,bookmarks=false]{hyperref}
\newcommand{\ie}{i.e.\ }
\newcommand{\cf}{cf.\ }
\newcommand{\resp}{respectively\ }
\newcommand{\Subsection}{\subsection}
\providecommand{\nobreakdash}{\nobreak\hbox{-}}
\setlength{\emergencystretch}{2em}
\multlinegap0pt
\usepackage{bm}
\usepackage{bbm}
\usepackage{dsfont}
\usepackage{xspace}
\usepackage{cancel}
\usepackage{mathtools}
\usepackage{amsthm}
\makeatletter
\@ifundefined{theorem}{\newtheorem{theorem}{Theorem}}{}
\@ifundefined{defn}{\newtheorem{defn}[theorem]{Definition}}{}
\@ifundefined{proposition}{\newtheorem{proposition}[theorem]{Proposition}}{}
\makeatother
\newcommand{\ZZ}{\mathbb{Z}}
\newcommand{\etale}{\textup{\'et}}
\title[Trito-non-ordinary Iwasawa theory of diagonal cycles]{Trito-non-ordinary Iwasawa theory of diagonal cycles}
\author{Ra\'ul Alonso, K\^az{\i}m B\"uy\"ukboduk, Antonio Cauchi,, Antonio Lei}
\date{Source: arXiv:2511.11511v1 (CC BY 4.0)}
\begin{document}
\maketitle

\noindent\emph{Source and licence.} Trito-non-ordinary Iwasawa theory of diagonal cycles. Version arXiv:2511.11511v1 (CC BY 4.0), distributed under the Creative Commons Attribution 4.0 International licence, as recorded in the arXiv licence metadata for this exact version and shown on the abstract page https://arxiv.org/abs/2511.11511v1. Full rights notice: licenses/arxiv-2511.11511-CC-BY-4.0.txt.\par\smallskip

\noindent\textit{Editorial scope.} Counted unit: the Proposition whose source label is \texttt{prop\_2025\_03\_08\_1713} in the article (source0.tex lines 1130--1134, source label \texttt{prop\_2025\_03\_08\_1713}), delivered together with the article's own complete original proof of it (source0.tex lines 1135--1261, one \texttt{proof} environment of 127 lines). No mathematics is written, repaired or re-derived: statement and proof are the article's own text line for line. Required context retained uncounted: prop\_2025\_05\_15\_0828 (lines 729--733); def\_2025\_07\_03\_1206 (lines 842--845). Every definition, display and label of those lines is retained unchanged. Omitted from this package and not redistributed: 1--728, 734--841, 846--1129, 1262--3279. Nothing inside a retained excerpt is omitted or altered: every retained body is delivered exactly as the article's lines are written, so no omission receipt of any kind accompanies this package. Citations: the retained excerpts carry no citation command at all, so no bibliography entry of the article is transcribed and no reference is redistributed. Reference adaptation: none was needed and none was made; every reference inside the delivered unit resolves inside it, so no unresolved reference or citation is produced. Printed slips kept literal: no printed word, symbol, hypothesis or number of the counted statement or of its proof was altered. Layout and class: the article's own class is \texttt{amsart} with packages including dsfont, xcolor, inputenc, fontenc, tipa, helvet, color, mathrsfs, stmaryrd, amsmath, amsxtra, amsthm, amssymb, xr; the wrapper is the lane's pinned \texttt{amsart} editorial wrapper at 10pt on A4, which restates the theorem-like environments the retained text needs and adds the article's own macro definitions that the retained text needs (\texttt{\textbackslash ZZ}, \texttt{\textbackslash etale}), with extra wrapper packages (bm, bbm, dsfont, xspace, cancel, mathtools). The delivered numbering is the unit's own. Assets: no retained span contains a figure environment, a graphics-inclusion command, a PDF-page inclusion, a table or a TikZ picture; no image file is copied into this package, the package declares no asset dependency and no image file is redistributed with it. Licence: version arXiv:2511.11511v1 of this article is distributed under the Creative Commons Attribution 4.0 International licence, as recorded in the arXiv licence metadata for this exact version and shown on the abstract page https://arxiv.org/abs/2511.11511v1. A full rights notice is delivered at licenses/arxiv-2511.11511-CC-BY-4.0.txt. Characters: every retained excerpt is pure ASCII, so the delivered engine, XeLaTeX with its default Latin Modern text font, reports no missing character.
\begin{proposition}
\label{prop_2025_05_15_0828}
    Let $T$, resp.  $\widetilde{T}$, be the Hecke correspondences in $\mathrm{T}_{B_E}(Np^{\mathbf{n}})$, resp. $\mathrm{T}_{B_E,{\rm spl}}(Np^{\mathbf{n}})$, associated with $g \in \mathcal{O}_{B_E,f}^\times$. Then
     \[ {\rm b}_{\mathbf{n}}^* \circ T = \widetilde{T}  \circ {\rm b}_{\mathbf{n}}^*, \text{ and } {\rm b}_{\mathbf{n},\ast} \circ  \widetilde{T} = T \circ {\rm b}_{\mathbf{n},\ast}.\]
\end{proposition}

\begin{defn}\label{def_2025_07_03_1206}
    \item[i)] For $n \in \mathbb N$, we define $\Delta_{\mathcal{U}_{(n)}} := \iota(\eta_p^n)_\ast \circ \iota_{n,\ast}^u(\mathbf{1}_{n}) \in {\rm CH}^2(\mathbf{X}_{\mathcal{U}_{(n)}})$.        
    \item[ii)]    For every ${\mathbf{n} } = (n_1,n_2,n_3) \in \mathbb N^3$, we define $\Delta_{\mathcal{U}_{\mathbf{n}}} : = {\rm pr}^{(n)}_{\mathbf{n}} \Delta_{\mathcal{U}_{(n)}} $, where $n = \max \{ n_1,n_2,n_3 \}$.
\end{defn}

\begin{proposition}
        \label{prop_2025_03_08_1713}
Let $n\geq 2$ be an integer. For any integer $d\equiv 1 \mod p^{n-1}$, we have
    $$\langle 1,1,d\rangle\,\widetilde{\Delta}_{n}^{\etale}(f, g)=\widetilde{\Delta}_{n}^{\etale}(f,g)\,.$$
\end{proposition}

\begin{proof}
    %\item[i)] 
    We recall that 
$$\widetilde{\Delta}_{n}^{\etale}(f,g)=({\rm pr}_f,\,{\rm pr}_g\,,\,{\rm id})\circ ({\rm id},\pi_2,{\rm id})_*\circ {\rm pr}^{(n)}_{1,1,n}  \,\Delta_{(n)}^{\ddagger}\,.$$
Hence, by obvious commutation relations (including $\pi_{2,*}\circ \pi_{1,*}=\pi_{1,*}\circ\pi_{2,*}$), we have
\begin{align}
\label{eqn_2025_03_09_1004}
\begin{aligned}
    \langle 1,1,d^{-1}\rangle\,\widetilde{\Delta}_{n}^{\etale}(f,g) =({\rm pr}_f,\,{\rm pr}_g\,,\,{\rm id})\circ ({\rm id},\pi_1,{\rm id})_*\circ \,{\rm pr}^{n-1,n-1,n}_{1,1,n}\circ \langle 1,1,d^{-1}\rangle \circ (\pi_1,\pi_2,{\rm id})_*\,\Delta_{(n)}^{\ddagger}\,.
\end{aligned}
\end{align}
Recall also that 
$$\Delta_{(n)}^{\ddagger}:=C_q^{-1}(\mathrm{U}_{p}',{\rm id},\mathrm{U}_{p}')^{-n}\Delta_{(n)}^{\etale}\,,$$
so that, since $({\rm U}_p', {\rm id},{\rm U}_p')$ commutes with diamond action and $(\pi_1,\pi_2,{\rm id})_*$ (note that $n\geq 2$),
\begin{equation}
\label{eqn_2025_03_09_1005}
\langle 1,1,d^{-1}\rangle \circ (\pi_1,\pi_2,{\rm id})_* \,\Delta_{(n)}^{\ddagger}=C_q^{-1}(\mathrm{U}_{p}',{\rm id},\mathrm{U}_{p}')^{-n}\langle 1,1,d^{-1}\rangle  \circ (\pi_1,\pi_2,{\rm id})_*\,\Delta_{(n)}^{\etale}\,.
\end{equation}
Combining \eqref{eqn_2025_03_09_1004} and \eqref{eqn_2025_03_09_1005}, it therefore suffices (in order to check the claimed equality in the first part of our proposition) to prove that
\begin{equation}
\label{eqn_2025_03_22_1357}
     \langle 1,1,d^{-1}\rangle \circ (\pi_1,\pi_2,{\rm id})_*\,\Delta_{(n)}^{\etale}=(\pi_1,\pi_2,{\rm id})_*\,\Delta_{(n)}^{\etale}\,.
\end{equation}
Furthermore, as we have (by definition)
\begin{equation}
\label{eqn_2025_03_22_1414}
    \Delta_{(n)}^{\etale}:= 
%\texttt{K\"u}\circ 
{\rm cl}_1 \circ (e_{\mathrm{U}_{p}'}, \mathrm{T}_{q}'-(q+1), e_{\mathrm{U}_{p}'}) \cdot\,\Delta_{(n)} %\circ \iota(\eta_p^n)_\star \circ \iota_{n,\star}^u(\mathbf{1}_{n})\,, 
\end{equation}
and $\langle 1,1,d^{-1}\rangle \circ (\pi_1,\pi_2,{\rm id})_*$ commutes with the
%K\"unneth isomorphism $\texttt{K\"u}$, 
cycle class map ${\rm cl}_1$ as well as the ordinary projector (we use the fact that $n\geq 2$ once again), it suffices to show that
\begin{equation}
\label{eqn_2025_07_02_1642}
\langle 1,1,d^{-1}\rangle \circ (\pi_1,\pi_2,{\rm id})_*\,\Delta_{(n)}=(\pi_1,\pi_2,{\rm id})_*\,\Delta_{(n)}\,. 
\end{equation}
%Since $\Delta_{(n)} = {\rm b}_{(n),*}\Delta_{\mathcal{U}_{(n)}}$, the sought-after equality \eqref{eqn_2025_07_02_1642} follows (thanks to Proposition \ref{prop_2025_05_15_0828}) once we show 
%\begin{equation}
%\label{eqn_2025_07_02_1647} 
%\langle 1,1,d^{-1}\rangle \circ (\pi_1,\pi_2,{\rm id})_*\,\Delta_{\mathcal{U}_{(n)}}= (\pi_1,\pi_2,{\rm id})_* \circ\,\langle 1,1,d^{-1}\rangle\,\Delta_{\mathcal{U}_{(n)}}=(\pi_1,\pi_2,{\rm id})_*\,\Delta_{\mathcal{U}_{(n)}}\,, 
%\end{equation}
%where the first equality is clear. In \eqref{eqn_2025_07_02_1647}, we have abusively denoted by $(\pi_1,\pi_2,{\rm id})$ the morphism $\mathbf{X}_{\mathcal{U}_{(n)}} \to \mathbf{X}_{\mathcal{U}_{n-1,n-1,n}}$ given as the compositum 
%\[\mathbf{X}_{\mathcal{U}_{(n)}} \xrightarrow[]{\mathcal{T}_{{}^2\eta_p^{-1}}} \mathbf{X}_{{}^2\eta_p \,\mathcal{U}_{(n)}\,{}^2\eta_p^{-1} } \xrightarrow{\,\,{\rm pr}\,\,} \mathbf{X}_{\mathcal{U}_{n-1,n-1,n}}, \]
%where we recall that ${}^2\eta_p = \left({\rm id},\eta_p ,{\rm id}\right)$ and ${\rm pr}$ denotes the natural degeneracy map given by the inclusion ${}^2\eta_p\, \mathcal{U}_{(n)}\,{}^2\eta_p^{-1} \subseteq \mathcal{U}_{n-1,n-1,n}$.


%{(2025/07/02) We should take one more step back and move to non-split tower and the class $\Delta_{\mathcal{U}_{(n)}}$, and recast this sought after identity in terms of that.}

%{(2025/07/03)}


Since $d\equiv 1 \pmod{p^{n-1}}$, the elements $\left ( \begin{smallmatrix}
    d & \\ & 1
\end{smallmatrix}\right), \left ( \begin{smallmatrix}
    d & (d-1)p^{1-n}\\ & 1
\end{smallmatrix}\right) \in \mathrm{GL}_2(\ZZ_p)$ belong to $U_{n-1}$, and hence
\begin{align*}
    \langle 1,1,d^{-1}\rangle \circ (\pi_1,\pi_2,{\rm id})_*\,\Delta_{(n)} &= \left (\left ( \begin{smallmatrix}
    d & \\ & 1
\end{smallmatrix}\right), \left ( \begin{smallmatrix}
    d & (d-1)p^{1-n}\\ & 1
\end{smallmatrix}\right), \left ( \begin{smallmatrix}
    1 & \\ & d
\end{smallmatrix}\right)\right)_* (\pi_1,\pi_2,{\rm id})_*\,\Delta_{(n)} \\ &= (\pi_1,\pi_2,{\rm id})_* \left (\left ( \begin{smallmatrix}
    d & \\ & 1
\end{smallmatrix}\right), \left ( \begin{smallmatrix}
    d & (d-1)p^{-n}\\ & 1
\end{smallmatrix}\right), \left ( \begin{smallmatrix}
    1 & \\ & d
\end{smallmatrix}\right)\right)_* \,\Delta_{(n)}.
\end{align*}
Moreover, as $\Delta_{(n)} = {\rm b}_{(n),*}\Delta_{\mathcal{U}_{(n)}}$, the sought-after equality \eqref{eqn_2025_07_02_1642} follows (thanks to Proposition \ref{prop_2025_05_15_0828}) once we show that
\begin{equation}
\label{eqn_2025_07_03_1204} 
\left (\left ( \begin{smallmatrix}
    d & \\ & 1
\end{smallmatrix}\right), \left ( \begin{smallmatrix}
    d & (d-1)p^{-n}\\ & 1
\end{smallmatrix}\right), \left ( \begin{smallmatrix}
    1 & \\ & d
\end{smallmatrix}\right)\right)_* \,\Delta_{\mathcal{U}_{(n)}}=\,\Delta_{\mathcal{U}_{(n)}}\,. 
\end{equation}
By Definition \ref{def_2025_07_03_1206}, we have
\begin{align*}
    \left (\left ( \begin{smallmatrix}
    d & \\ & 1
\end{smallmatrix}\right), \left ( \begin{smallmatrix}
    d & (d-1)p^{-n}\\ & 1
\end{smallmatrix}\right), \left ( \begin{smallmatrix}
    1 & \\ & d
\end{smallmatrix}\right)\right)_* \,\Delta_{\mathcal{U}_{(n)}} & = \left (\left ( \begin{smallmatrix}
    d & \\ & 1
\end{smallmatrix}\right), \left ( \begin{smallmatrix}
    d & (d-1)p^{-n}\\ & 1
\end{smallmatrix}\right), \left ( \begin{smallmatrix}
    1 & \\ & d
\end{smallmatrix}\right)\right)_* \circ  \iota(\eta_p^n)_\ast \circ \iota_{n,\ast}^u(\mathbf{1}_{n}) \\ 
& =\iota(\eta_p^n)_\ast \circ \left (\left ( \begin{smallmatrix}
    d & \\ & 1
\end{smallmatrix}\right), \left ( \begin{smallmatrix}
    d & d-1\\ & 1
\end{smallmatrix}\right), \left ( \begin{smallmatrix}
    1 & \\ & d
\end{smallmatrix}\right)\right)_* \circ  \iota_{n,\ast}^u(\mathbf{1}_{n}) \\ 
& =\iota(\eta_p^n)_\ast \circ  u_*  \circ \iota \left (\left ( \begin{smallmatrix}
    d & \\ & 1
\end{smallmatrix}\right)\right)_* \circ \iota_{U_{Z,n}\,,\, u\, \mathcal{U}^{(n)} \,u^{-1},\ast}(\mathbf{1}_{n}) 
\\ 
& =\iota(\eta_p^n)_\ast \circ  u_*   \circ \iota_{U_{Z,n}\,,\, u\, \mathcal{U}^{(n)}\, u^{-1},\ast}(\left ( \begin{smallmatrix}
    d & \\ & 1
\end{smallmatrix}\right)_*\mathbf{1}_{n})\,. 
\end{align*}
Recall that $$\mathbf{1}_{n}= \sum_{\mathcal{C} \in \pi_0(X_{Z,n})} [\mathcal{C}] \in {\rm CH}^0(X_{Z,n})$$  
is the fundamental cycle of $X_{Z,n}$.
Since $d\equiv 1 \mod p^{n-1}$ and $n>1$, it follows that $d$ is a square mod $p^n$. As a result, the action of $\left ( \begin{smallmatrix}
    d & \\ & 1
\end{smallmatrix}\right)_*$ preserves the set of connected components of $X_{Z,n}$. Then the equality \eqref{eqn_2025_07_03_1204} follows from the fact that 
$$\left ( \begin{smallmatrix}
    d & \\ & 1
\end{smallmatrix}\right)_*[\mathcal{C}]=\left ( \begin{smallmatrix}
    d^{-1} & \\ & 1
\end{smallmatrix}\right)^*[\mathcal{C}]=[\mathcal{C}]\,,$$
where the last equality is because the fundamental cycle is compatible with (flat) base change and $\left ( \begin{smallmatrix}
    d^{-1} & \\ & 1
\end{smallmatrix}\right)^*$ is an isomorphism (in particular, its degree equals $1$). This concludes the proof.
\end{proof}

\end{document}
